\section{Introduction}\label{sec:introduction}

Modern agent systems compose identities, tools, model providers, delegated
permissions, jurisdictional claims, policy decisions, and evidence artifacts.
In many deployments those layers live in different places: prompts describe
intent, configuration names providers, application code enforces some checks,
logs record partial outcomes, and documents describe governance after the
fact.  This separation is operationally convenient, but it makes a systems
question hard to answer: what can be represented, restricted, and evidenced
together before an agent operation reaches an executable provider?

\Argorix answers this question as a compiled governance runtime rather than as
a certification framework.  It compiles agent declarations, typed messages,
local Agent Passport metadata, provider contracts, policies, runtime profiles,
and evidence requirements into bytecode that is executed by a fail-closed VM
\citep{argorixlang2026}.  The VM separates a declarative provider contract from
an executable provider adapter.  A contract may describe an external provider;
it does not become callable unless an executable adapter is registered.  In the
artifact studied here, \code{simulated} is the only executable provider.

The central thesis is claim-bounded: \Argorix demonstrates a language/runtime
boundary that can compile governance metadata, deny unsupported provider
execution, and generate internally verifiable local evidence.  It does not
turn declarations into external truth.  Throughout the paper we enforce five
separations:
\[
\begin{array}{rcl}
\mathrm{declaration} & \ne & \mathrm{verification},\\
\mathrm{digest\ integrity} & \ne & \mathrm{policy\ approval},\\
\mathrm{local\ passport\ metadata} & \ne & \mathrm{authenticated\ identity},\\
\mathrm{provider\ contract} & \ne & \mathrm{executable\ adapter},\\
\mathrm{EvidenceBundle} & \ne & \mathrm{external\ attestation}.
\end{array}
\]
These separations are not caveats appended at the end; they define the method.

Related systems occupy adjacent points in the design space.  MCP and A2A
standardize agent interaction surfaces \citep{mcp_spec_2025,a2a_spec}.  NANDA
and AgentFacts motivate open naming and discovery for agents
\citep{projectnanda2026,raskar2025beyonddns}.  DID and VC standards define
identifier and credential data models
\citep{w3c_did_core_2022,w3c_vc_data_model_2025}.  OPA and Cedar demonstrate
policy-as-code approaches \citep{opa_docs,cedar_policy}; SPIFFE/SPIRE address
workload identity \citep{spiffe_overview,spire_overview}; in-toto and SLSA
address provenance and supply-chain evidence
\citep{torresarias2019intoto,slsa_spec}.  \Argorix does not replace these
systems.  Its contribution is to compile local governance constructs and
evidence boundaries in one agent runtime while marking unresolved assurance
properties as unresolved.

The paper asks five explicit research questions:
\begin{description}
\item[RQ1.] How can governance, sovereign identity, and evidence requirements
be represented as compiled language constructs?
\item[RQ2.] How does the \Argorix pipeline constrain a chatbot request before a
provider operation can be planned?
\item[RQ3.] What evidence does the runtime generate, and how can its internal
consistency be checked offline?
\item[RQ4.] What do the existing runtime sessions demonstrate, and which
controlled experiments remain required?
\item[RQ5.] How can Agent Passports, ATrust, DCP-AI, and future sovereign
discovery fit into an open agentic-web architecture?
\end{description}

Our contributions are:
\begin{enumerate}
\item \textbf{Compiled governance constructs:} a language architecture that
represents agents, passports, policies, runtime profiles, provider contracts,
and evidence requirements in one validated pipeline.
\item \textbf{Fail-closed provider boundary:} a split between declarative
provider contracts and executable adapters, with unsupported external-provider
execution denied before provider planning.
\item \textbf{Typed policy lattice:} a v0.2 controlled-matrix policy model
separating pass, deny, review, unknown-rule configuration errors, and malformed
policy objects instead of collapsing unknown rules into ordinary violations.
\item \textbf{Offline evidence boundary:} an EvidenceBundle model that checks
bytecode, trace, security-report, and ledger digest consistency while
explicitly excluding source integrity, authorship, timestamp, and side-effect
absence from the current claim.
\item \textbf{Controlled evaluation agenda:} a snapshot analysis plus generated
controlled-case matrix, ablations, baselines, and claim-boundary taxonomy.
\end{enumerate}

To prevent overreading, we classify every claim into one of four categories from
the outset, summarized in \cref{tab:claim-boundaries} and applied uniformly
thereafter.  \emph{Implemented} mechanisms include the compiler, the fail-closed
VM path, the simulated-only provider boundary, traces, the SecurityReport, the
EvidenceBundle, and semantic digest verification.  \emph{Declarative} metadata
includes passports, \code{ans_name}, \ATrust maps, MCP and A2A bridge
declarations, and external-provider contracts.  \emph{Proposed} architecture
includes Sovereign DNS/ANS resolution, a DID/VC resolver, signed bundles,
source-digest binding, and live federation.  We explicitly do \emph{not} claim
real identity authentication, credential verification, cryptographic
handshakes, blockchain immutability, post-quantum security, legal-compliance
certification, universal prompt-injection prevention, or production isolation.

\begin{figure}[htbp]
\centering
\includegraphics[width=\linewidth]{architecture.pdf}
\caption{The governed execution and sovereign trust stack. Solid paths denote
implemented compilation and runtime behavior; metadata layers do not imply
external verification.}
\label{fig:architecture}
\end{figure}

\Cref{fig:architecture} summarizes the stack.  The empirical result is
deliberately narrower than an assurance claim: we observe deterministic
structure, bounded execution, and verifiable artifact relationships in one
repeated snapshot.  We neither conduct a controlled attack experiment nor infer
production security from successful digest verification.

Where a figure and table address the same topic, they are deliberately paired
rather than duplicative: the figure communicates topology, flow, or boundary
placement, while the table provides exact categorical or quantitative
comparison.  Captions and prose identify the distinct evidentiary role of each.
